\documentclass[10pt,a4paper]{amsart}
\usepackage[margin=19mm]{geometry}
\usepackage{amsmath,amssymb,mathtools}
\usepackage[scr=rsfs,cal=euler]{mathalfa}
\usepackage{xfrac}
\usepackage[unicode,hidelinks,bookmarks=false]{hyperref}
\newcommand{\ie}{i.e.\ }
\newcommand{\cf}{cf.\ }
\newcommand{\resp}{respectively\ }
\newcommand{\Subsection}{\subsection}
\providecommand{\nobreakdash}{\nobreak\hbox{-}}
\setlength{\emergencystretch}{2em}
\multlinegap0pt
\usepackage{bm}
\usepackage{bbm}
\usepackage{dsfont}
\usepackage{xspace}
\usepackage{cancel}
\usepackage{mathtools}
\usepackage{amsthm}
\makeatletter
\@ifundefined{theorem}{\newtheorem{theorem}{Theorem}}{}
\@ifundefined{lemma}{\newtheorem{lemma}[theorem]{Lemma}}{}
\makeatother
\title[$G\_2$ and the Maximally Symmetric (3, 8) Distribution with 6]{$G\_2$ and the Maximally Symmetric (3, 8) Distribution with 6-Dimensional Square}
\author{Nicklas Day, Boris Doubrov,, Igor Zelenko}
\date{Source: arXiv:2605.25910v3 (CC BY 4.0)}
\begin{document}
\maketitle

\noindent\emph{Source and licence.} $G_2$ and the Maximally Symmetric (3, 8) Distribution with 6-Dimensional Square. Version arXiv:2605.25910v3 (CC BY 4.0), distributed under the Creative Commons Attribution 4.0 International licence, as recorded in the arXiv licence metadata for this exact version and shown on the abstract page https://arxiv.org/abs/2605.25910v3. Full rights notice: licenses/arxiv-2605.25910-CC-BY-4.0.txt.\par\smallskip

\noindent\textit{Editorial scope.} Counted unit: the Lemma whose source label is \texttt{lemma: Delta' growth} in the article (source0.tex lines 1647--1650, source label \texttt{lemma: Delta' growth}), delivered together with the article's own complete original proof of it (source0.tex lines 1651--1715, one \texttt{proof} environment of 65 lines). No mathematics is written, repaired or re-derived: statement and proof are the article's own text line for line. Required context retained uncounted: omega\_1\_p (lines 1625--1628). Every definition, display and label of those lines is retained unchanged. Omitted from this package and not redistributed: 1--1624, 1629--1646, 1716--2359. Nothing inside a retained excerpt is omitted or altered: every retained body is delivered exactly as the article's lines are written, so no omission receipt of any kind accompanies this package. Citations: the retained excerpts carry no citation command at all, so no bibliography entry of the article is transcribed and no reference is redistributed. Reference adaptation: none was needed and none was made; every reference inside the delivered unit resolves inside it, so no unresolved reference or citation is produced. Printed slips kept literal: no printed word, symbol, hypothesis or number of the counted statement or of its proof was altered. Layout and class: the article's own class is \texttt{amsart} with packages including graphicx, standalone, float, amsthm, amssymb, amsmath, hyperref, dsfont, enumitem, tcolorbox, esint, mathtools, contour, appendix; the wrapper is the lane's pinned \texttt{amsart} editorial wrapper at 10pt on A4, which restates the theorem-like environments the retained text needs and adds the article's own macro definitions that the retained text needs (none), with extra wrapper packages (bm, bbm, dsfont, xspace, cancel, mathtools). The delivered numbering is the unit's own. Assets: no retained span contains a figure environment, a graphics-inclusion command, a PDF-page inclusion, a table or a TikZ picture; no image file is copied into this package, the package declares no asset dependency and no image file is redistributed with it. Licence: version arXiv:2605.25910v3 of this article is distributed under the Creative Commons Attribution 4.0 International licence, as recorded in the arXiv licence metadata for this exact version and shown on the abstract page https://arxiv.org/abs/2605.25910v3. A full rights notice is delivered at licenses/arxiv-2605.25910-CC-BY-4.0.txt. Characters: every retained excerpt is pure ASCII, so the delivered engine, XeLaTeX with its default Latin Modern text font, reports no missing character.
    \begin{gather}
        \label{omega_1_p}
        (\omega_1)_p = dx_1^5\cdot p_0 + (dx_1^2+dx_1^7)\cdot z_1 + (dx_1^3-dx_1^6)\cdot z_2
    \end{gather}

\begin{lemma}
    \label{lemma: Delta' growth}
    The square $(\Delta')^2$ of $\Delta'$ has rank $10$, and $(\Delta')^3 = TTK$.
\end{lemma}

\begin{proof}
    As above, fix $p_0 = i + \ell$, $z_1 = k+\ell j$, and $z_2 = -j +\ell k$, so that $\Delta_{p_0} = \mathrm{span}\{p_0,z_1,z_2\}$. Because $\Delta'$ is homogeneous, it suffices to show that $\dim\big((I_{\Delta'}^{(1)})_{p}\big) = 2$, where $p= (p_0,0)\in TK$. 
    
    Fix a vector $u\in \mathbb{R}^7$ and a vector field $v\in \Gamma(TK)$ to obtain an arbitrary differential 1-form
    \begin{equation}
        \label{def varphi}
        \varphi =Q(u,\pi^*\omega_0) +  Q(v, \omega_1) \in I_{\Delta'}
    \end{equation}
    and suppose that $d\varphi|_{\Delta'}=0$ so that $\varphi\in I_{\Delta'}^{(1)}$. 

    Firstly, we show that $Q(v,\omega_1)_{p} = 0$. Fix $z_3 = \ell i$, and observe that 
    \begin{gather}
        \label{z3 cross conditions}
        p_0\times z_3 = p_0, \quad z_1\times z_3 = -z_1,\quad \text{and}\quad z_2\times z_3 = -z_2.
    \end{gather}
    We can compute
    \begin{equation}
        \label{dN}
        d\big(Q(v,\omega_1)\big) = Q(dv,\omega_1) + Q(v,d\omega_1)
    \end{equation}
    where $dv$ is the differential of the function $v:TK\to \mathbb{R}^7$. Consider the vectors $(z_1,0),(z_2,0),$ and $(0,z_3)\in \Delta'_p$. For $a=1,2$, we can use the definitions of $\omega_0$ and $\omega_1$ to compute
    \begin{gather}
        0=d\varphi\big((z_a,0)\wedge(0,z_3)\big) = d\big(Q(v,\omega_1)\big)\big((z_a,0)\wedge(0,z_3)\big) 
        \\
        = Q\Big(dv(z_a), \omega_1\big((0,z_3)\big)\Big) + Q\Big(v,d\omega_1\big((z_a,0)\wedge (0,z_3)\big)\Big)
        \\
        = Q\big(dv(z_a),x_0\times z_3\big) + Q(v, 2\cdot z_a\times z_3)
        = Q\big(dv(z_a),x_0\big) - 2\cdot Q(v, z_a)
    \end{gather}
    where the last equality is by \eqref{z3 cross conditions}. Because $Q(x_0,v)=0$, we have that
    \begin{equation}
    \label{dv cond}
        Q(v,dx_0) + Q(dv,x_0)=0,
    \end{equation} so we can continue the above equalities
    \begin{gather}
        Q\big(dv(z_a),x_0\big) - 2\cdot Q(v, z_a) = - 3\cdot Q(v,z_a),
    \end{gather}
    and therefore $Q(v,z_a) = 0$ for $i=1,2$. Since $\Delta_{p_0} = \mathrm{span}\{p_0,z_1,z_2\}$ and $v_{p_0}\in T_{p_0}K$, this implies that $v_{p_0}\in \Delta^\perp_{p_0}$. Because \eqref{omega_1_p} shows $(\omega_1)_p$ takes values in $\Delta_{p_0}$, we have now shown that $Q(v,\omega_1)_{p}=0$ and hence that $\varphi = Q(u,\pi^*\omega_0)$.

    We can now see that
    \begin{equation}
        d\varphi = Q(u,d\omega_0) = Q(u,2dx_0\times dx_0).
    \end{equation}
    Again restricting to the point $p$, we have that 
    \begin{gather}
        d\varphi_{p}\big((p_0,0),(z_1,0)) = Q(u,2p_0\times z_1) = 0
        \\
        d\varphi_{p}\big((p_0,0),(z_2,0)) = Q(u,2p_0\times z_2) = 0
        \\
        d\varphi_{p}\big((z_1,0),(z_2,0)) = Q(u,2z_1\times z_2) = Q(u,4p_0).
    \end{gather}
    Therefore, $\varphi_{p}\in (I_{\Delta'}^{(1)})_p$ if and only if $u_{p_0}\in p_0^\perp = T_{p_0}K$. Since 
    \[
        Q\big(u_{p_0},\pi^*(\omega_0)_{p_0}\big) \equiv 0\Longleftrightarrow u_{p_0}\in \Delta^\perp_{p_0}\subseteq T_{p_0}K,
    \]
    we have that
    \begin{equation}
        \dim\big(I_{\Delta'}^{(1)}\big)_{p_0} = \dim(T_{p_0}K) - \dim(\Delta_{p_0}^\perp) = 6-4 = 2.
    \end{equation}
    This demonstrates that $(\Delta')^2 = \pi^*\big(\Delta^2\big)$ has rank $10$. To see that $(\Delta')^3 = TTK$, observe that
    \begin{equation}
        \pi_*\big((\Delta')^{3}\big) = \pi_*\big(\pi^*\big(\Delta^2\big) + [\Delta',\pi^*\big(\Delta^2\big)]\big) = \Delta^2 + [\Delta, \Delta^2] = \Delta^3 = TK
    \end{equation}
    and that since $(\Delta')^2 = \pi^*(\Delta^2)$ contains the vertical subbundle for the projection $TTK\to TK$, so does $(\Delta')^3$.
\end{proof}

\end{document}
